\section{Экспериментальная часть}

\subsection{Постановка задачи}

Эмпирическая часть работы направлена на проверку того, как разные группы признаков влияют на прогноз направления движения цены. Задача формулируется как бинарная классификация: для каждого дня $t$ необходимо определить, будет ли цена закрытия следующего торгового дня выше цены закрытия текущего дня.

Целевая переменная задаётся следующим образом:

\[
y_t = \mathbf{1}\{Close_{t+1} > Close_t\}.
\]

Если цена закрытия на следующий торговый день выше цены закрытия текущего дня, наблюдению присваивается класс $1$; если цена не выросла, присваивается класс $0$. Такая постановка позволяет оценивать не точное значение цены, а направление движения актива.

Главный исследовательский вопрос экспериментальной части состоит в следующем: дают ли признаки, извлечённые из финансовых текстов с помощью FinBERT, самостоятельный сигнал и добавочную предсказательную силу относительно рыночных признаков.

\subsection{Данные}

В качестве объекта эксперимента выбран актив AAPL. Данные имеют дневную частоту, а горизонт прогноза составляет один торговый день. Итоговый набор данных основного запуска содержит 1578 наблюдений.

Для контроля утечки данных используется хронологическое разбиение по дате принятия решения. Одна и та же дата не может попадать одновременно в train, validation и test. Основной запуск RUN 03 использует split \texttt{chronological\_row\_balanced\_by\_decision\_date}.

\begin{table}[h]
\centering
\begin{tabular}{lrrl}
\toprule
Подвыборка & Наблюдения & Уникальные даты & Период \\
\midrule
Train & 1104 & 52 & 2020-04-01 --- 2020-06-15 \\
Validation & 252 & 10 & 2020-06-16 --- 2020-06-29 \\
Test & 222 & 55 & 2020-06-30 --- 2020-09-30 \\
\bottomrule
\end{tabular}
\caption{Хронологическое разбиение данных в RUN 03}
\end{table}

Такое разбиение сохраняет временную структуру данных и одновременно оставляет test subset достаточного размера для первичного сравнения групп признаков. Текстовые признаки используются с лагом: для прогноза после дня $t$ учитываются тексты не позднее $t-1$. Это правило снижает риск попадания в модель информации, которая фактически появилась после прогнозируемого события.
